%% Courbe de chauffage de 1,0 kg d'eau, de -20 °C à 120 °C, à puissance constante.
%% La puissance étant constante, la LONGUEUR de chaque portion en abscisse est
%% proportionnelle à l'énergie reçue (bilan : 42 + 334 + 420 + 2260 + 40 = 3096 kJ).
%% Échelle des abscisses : 3096 kJ -> 6,8 cm, soit 0,0021964 cm par kJ
%%   x1 = 0,092 | x2 = 0,826 | x3 = 1,748 | x4 = 6,712 | x5 = 6,800
%% Échelle des ordonnées : 1 °C -> 0,0215 cm, origine à -20 °C.
\documentclass[tikz,border=4pt]{standalone}
\usepackage{liaisons}
\begin{document}
\begin{tikzpicture}[x=1cm,y=1cm,
    axe/.style={-{Stealth[length=5pt]}, line width=0.6pt},
    courbe/.style={solideB, line width=1.6pt, line join=round, line cap=round},
    grad/.style={font=\scriptsize, inner sep=2pt},
    aide/.style={black!25, line width=0.4pt, dotted},
    repere/.style={line width=0.5pt},
    rappel/.style={line width=0.4pt, black!45},
    formule/.style={font=\scriptsize, anchor=west, inner sep=1pt}]

%% ---------- Macros locales ------------------------------------------------------
\newcommand\yt[1]{{(#1+20)*0.0215}}                % ordonnée d'une température en °C
\newcommand\pastille[3]{%                            % #1,#2 = position, #3 = numéro
  \fill[solideA] (#1,#2) circle (0.155);
  \node[font=\tiny\bfseries, text=white, inner sep=0pt] at (#1,#2) {#3};}

%% ---------- Axes et graduations -------------------------------------------------
\draw[aide] (0,\yt{0})   -- (7.05,\yt{0});
\draw[aide] (0,\yt{100}) -- (7.05,\yt{100});
\draw[aide] (0,\yt{120}) -- (7.05,\yt{120});
\draw[axe] (0,-0.10) -- (0,3.30);
\draw[axe] (-0.10,0)  -- (7.20,0) node[anchor=north east, inner sep=4pt, font=\small] {$t$ (s)};
\node[anchor=south west, inner sep=2pt, font=\small] at (0.05,3.32) {$\theta$ ($^\circ$C)};
\foreach \t/\lbl in {-20/{$-20$}, 0/{0}, 100/{100}, 120/{120}} {
  \draw[repere] (0,\yt{\t}) -- (-0.10,\yt{\t});
  \node[grad, anchor=east] at (-0.12,\yt{\t}) {\lbl};}

%% ---------- La courbe en cinq portions ------------------------------------------
\draw[courbe] (0,\yt{-20}) -- (0.092,\yt{0}) -- (0.826,\yt{0}) -- (1.748,\yt{100})
           -- (6.712,\yt{100}) -- (6.800,\yt{120});

%% ---------- Les deux paliers ----------------------------------------------------
\draw[rappel] (0.55,1.00) -- (0.47,\yt{4});
\node[font=\footnotesize, text=solideA, anchor=south] at (0.55,1.00) {Fusion};
\node[font=\footnotesize, text=solideA, anchor=south] at (4.23,\yt{104}) {Ébullition};
\draw[{Stealth[length=4pt]}-{Stealth[length=4pt]}, line width=0.6pt, black!60]
  (1.748,\yt{88}) -- (6.712,\yt{88});
\node[font=\scriptsize, fill=white, inner sep=1.5pt] at (4.23,\yt{88})
  {73 \% de l'énergie totale};

%% ---------- Numéros des portions ------------------------------------------------
\draw[rappel] (0.10,\yt{-11}) -- (0.30,\yt{-13});
\pastille{0.42}{\yt{-13}}{1}
\pastille{0.98}{\yt{0}}{2}
\pastille{1.287}{\yt{50}}{3}
\pastille{3.40}{\yt{100}}{4}
\draw[rappel] (6.76,\yt{110}) -- (6.45,\yt{113});
\pastille{6.30}{\yt{113}}{5}

%% ---------- Les cinq énergies ---------------------------------------------------
\newcommand\legformule[4]{\pastille{#1}{#2}{#3}\node[formule] at (#1+0.24,#2) {#4};}
\legformule{0.15}{-0.58}{1}{$Q_1 = m\,c_{\text{solide}}\,\Delta\theta$}
\legformule{3.75}{-0.58}{2}{$Q_2 = m\,L_{\text{fusion}}$}
\legformule{0.15}{-1.06}{3}{$Q_3 = m\,c_{\text{liquide}}\,\Delta\theta$}
\legformule{3.75}{-1.06}{4}{$Q_4 = m\,L_{\text{vaporisation}}$}
\legformule{0.15}{-1.54}{5}{$Q_5 = m\,c_{\text{gaz}}\,\Delta\theta$}
\node[font=\scriptsize, align=center, anchor=north] at (3.55,-1.90)
  {Puissance constante : chaque longueur est proportionnelle\\
   à l'énergie reçue (1,0 kg d'eau, $Q = 3\,096$ kJ).};
\end{tikzpicture}
\end{document}
